\section*{References}
\addcontentsline{toc}{section}{References}

\begin{enumerate}[leftmargin=1.7em,itemsep=0.5em]
\item \textbf{Sacred Scripture.} \emph{The Bible, Douay--Rheims Version},
Challoner revision, Project Gutenberg eBook~1581, registered verse-text
witness. \href{https://www.gutenberg.org/ebooks/1581}{Complete English text}.
Principal loci appear in the appendix on \hyperref[app:scripture]{scriptural loci}; the human
translation is public domain in the United States.

\item \textbf{Basil of Caesarea.} \emph{On the Holy Spirit}, XVI,
translated by Blomfield Jackson, \emph{Nicene and Post-Nicene Fathers},
second series, vol.~8 (1895).
\href{https://www.newadvent.org/fathers/3203.htm}{English text}.
The historical translation is public domain in the United States.

\item \textbf{Gregory Nazianzen.} Oration~38, nos.~9--10,
translated by Charles Gordon Browne and James Edward Swallow,
\emph{Nicene and Post-Nicene Fathers}, second series, vol.~7 (1894).
\href{https://www.newadvent.org/fathers/310238.htm}{English text}.

\item \textbf{Augustine.} \emph{The City of God}, chiefly XI--XII and XIV,
translated by Marcus Dods (1871; \emph{NPNF} reprint, first series,
vol.~2, 1887).
\href{https://www.newadvent.org/fathers/120111.htm}{Book XI}.
\emph{On the Trinity}, III, translated by Arthur West Haddan,
\emph{NPNF}, first series, vol.~3 (1887).
\href{https://www.newadvent.org/fathers/130103.htm}{Book III}.
Historical translations public domain in the United States.

\item \textbf{Dionysian corpus.} \emph{The Celestial Hierarchy}, in
\emph{The Works of Dionysius the Areopagite}, translated by John Parker,
vol.~II (London: James Parker and Co., 1899), pp.~1--66;
\emph{Patrologia Graeca}~3 for the chapter column addresses.
\href{https://www.tertullian.org/fathers/areopagite_13_heavenly_hierarchy.htm}{Complete
Parker transcription}. The traditional author attribution and historical
dating are explained in the body. The historical translation and Migne
edition are public domain in the United States.

\item \textbf{Dionysian dating and reception.} Benedict XVI,
general audience, 14 May 2008,
\href{https://www.vatican.va/content/benedict-xvi/en/audiences/2008/documents/hf_ben-xvi_aud_20080514.html}{official
English text}. Eric D. Perl, ``Pseudo-Dionysius the Areopagite,''
\emph{The Cambridge History of Philosophy in Late Antiquity},
\href{https://doi.org/10.1017/CHOL9780521194846.014}{publisher's public summary}.
Only the latter's public summary was consulted for the historical dating.

\item \textbf{Gregory the Great.} \emph{Homiliae in Evangelia},
II.34, Latin electronic witness, complete Wikisource export registered
17 September 2026.
\href{https://la.wikisource.org/wiki/Homiliarum_in_Evangelia/XXXIV}{Homily XXXIV}.
The underlying historical Latin and the modern export are distinct rights
objects; the registered export retains its own attribution and license.

\item \textbf{John Damascene.} \emph{Exposition of the Orthodox Faith},
II.3--4, translated by S.~D.~F. Salmond,
\emph{NPNF}, second series, vol.~9 (Edinburgh, 1898),
with the New Advent transcription as an additional witness.
\href{https://www.newadvent.org/fathers/33042.htm}{Book II}.

\item \textbf{Origen.} \emph{Commentary on Matthew}, XIII.26--29,
translated by John Patrick, \emph{Ante-Nicene Fathers}, vol.~9 (1896).
\href{https://www.newadvent.org/fathers/101613.htm}{Book XIII}.

\item \textbf{John Chrysostom.} \emph{Homilies on Matthew}, 59.4,
translated by George Prevost, revised by M.~B. Riddle,
\emph{NPNF}, first series, vol.~10 (1888).
\href{https://www.newadvent.org/fathers/200159.htm}{Homily 59}.

\item \textbf{Bonaventure.} \emph{Commentary on the Second Book of
Sentences}, distinction~3, part~1, article~1, questions~1--2,
\emph{Opera omnia}, vol.~2 (Quaracchi, 1885), pp.~89--91, 94.
\href{https://archive.org/details/doctorisseraphic02bona}{Facsimile}.
\emph{Breviloquium}, II.6--7,
\href{https://catholiclibrary.org/library/view?docId=/Medieval-OR/BonaventuraSBreviloquium.00000158.la.html}{Latin
web transcription}. The latter witness does not identify a controlling
printing; its use is confined to the clear Latin passages cited.

\item \textbf{Thomas Aquinas.} \emph{Summa theologiae}, I qq.~50--64,
106--114; III q.~8 a.~4 and q.~30, with additional loci cited locally.
Fathers of the English Dominican Province translation, second revised
edition (1920), \href{https://www.newadvent.org/summa/}{New Advent
transcription}; complete Prima Pars also registered from
\href{https://www.gutenberg.org/ebooks/17611}{Project Gutenberg 17611}.
\href{https://www.corpusthomisticum.org/sth1050.html}{Latin control for
I qq.~50--64}. \emph{De ente et essentia},
\href{https://www.corpusthomisticum.org/oee.html}{Baur--Koch text};
\emph{De veritate}, \href{https://www.corpusthomisticum.org/qdv08.html}{qq.~8--9};
\emph{De malo}, \href{https://www.corpusthomisticum.org/qdm16.html}{q.~16}.
Historical source wording retains its public-domain status; modern wrappers
and apparatus are separately identified in the source records.

\item \textbf{Fourth Lateran Council.} Constitution~1,
\emph{Firmiter credimus} (1215), DS~800--801;
\href{https://patristica.net/denzinger/enchiridion-symbolorum.html}{Latin
Denzinger witness}. \textbf{First Vatican Council.}
\emph{Dei Filius} (24 April 1870), ch.~1 and Creator canons~1--5;
\href{https://www.vatican.va/archive/hist_councils/i-vatican-council/documents/vat-i_const_18700424_dei-filius_la.html}{Holy
See Latin text}.

\item \textbf{Pius XII.} \emph{Humani generis} (12 August 1950), no.~26;
\href{https://www.vatican.va/content/pius-xii/en/encyclicals/documents/hf_p-xii_enc_12081950_humani-generis.html}{official
English witness}. \textbf{Second Vatican Council.} \emph{Lumen gentium}
(21 November 1964), 53, 56, 59, 66;
\href{https://www.vatican.va/archive/hist_councils/ii_vatican_council/documents/vat-ii_const_19641121_lumen-gentium_en.html}{official
English witness}.

\item \textbf{Catechism of the Catholic Church.} 328--336, 391--395;
\href{https://www.vatican.va/content/catechism/en/part_one/section_two/chapter_one/article_1/paragraph_5_heaven_and_earth.html}{Heaven
and earth};
\href{https://www.vatican.va/content/catechism/en/part_one/section_two/chapter_one/article_1/paragraph_7_the_fall.html}{the
fall}. Modern official text is cited and briefly paraphrased.

\item \textbf{Congregation for Divine Worship and the Discipline of the
Sacraments.} \emph{Directory on Popular Piety and the Liturgy}, decree
17 December 2001, nos.~213--217;
\href{https://www.vatican.va/roman_curia/congregations/ccdds/documents/rc_con_ccdds_doc_20020513_vers-direttorio_en.html}{official
English witness}. Source consulted for teaching and devotional direction;
no contemporary prayer translation is reproduced.

\item \textbf{Liturgical books.} \emph{Missale Romanum}, Vatican typical
edition (1962), angelic feasts at pp.~493, 670, 672, 695;
\href{https://media.churchmusicassociation.org/pdf/missale62.pdf}{facsimile}.
\emph{Ordo lectionum Missae}, second typical edition (1981),
nos.~647 and 650, pp.~308--309;
\href{https://archive.org/details/OLM1981}{facsimile}.
Metropolitan Cantor Institute, Byzantine Catholic Archeparchy of Pittsburgh,
\href{https://mci.archpitt.org/menaion/11-08.html}{Menaion, 8 November}.
Dates, ranks and appointments are reported from these identified witnesses;
modern liturgical text is not reproduced.
\end{enumerate}

Online witnesses were consulted on 29 September 2026 unless a retained
edition date is specified. The references identify the sources used;
their source-local qualifications concern the particular witnesses named.
